%\documentclass[a4paper,10pt]{article}
\documentclass[3p]{elsarticle}
%\usepackage[utf8]{inputenc}

%%%%%%%%%%%%%%%%%%%%%%%%%%%%%%%%%%%%%%%%%%%%%%%%%%%%%%%%%%%%%%%%%%%%%%%%
% Packages
\usepackage{graphicx}
\usepackage{amsmath,amssymb,amsfonts,amsthm}
\usepackage{xspace} 
\usepackage{hyperref}
%\usepackage{hyperref,todonotes}

\newcommand{\myleft}{left}
\newcommand{\myright}{right}

\graphicspath{{Figures/}}

%%%%%%%%%%%%%%%%%%%%%%%%%%%%%%%%%%%%%%%%%%%%%%%%%%%%%%%%%%%%%%%%%%%%%%%%

\title{Enriched finite elements and local rescaling for vibrations of axially inhomogeneous {Timoshenko} beams}

\author[irmar]{R.~Cornaggia\corref{cor1}}
\ead{remi.cornaggia@univ-rennes1.fr}

\author[irmar]{E.~Darrigrand}
\ead{eric.darrigrand@univ-rennes1.fr}

\author[irmar]{L.~Le~Marrec }
\ead{loic.lemarrec@univ-rennes1.fr}

\author[irmar]{F.~Mah\'e}
\ead{fabrice.mahe@univ-rennes1.fr}

\cortext[cor1]{Corresponding author \\ \emph{Email address}: remi.cornaggia@univ-rennes1.fr}

\address[irmar]{Universit\'e de Rennes, CNRS, IRMAR - UMR 6625, F-35000 Rennes, France}


%%%%%%%%%%%%%%%%%%%%%%%%%%%%%%%%%%%%%%%%%%%%%%%%%%%%%%%%%%%%%%%%%%%%%%%%
\begin{document}

%%%%%%%%%%%%%%%%%%%%%%%%%%%%%%%%%%%%%%%%%%%%%%%%%%%%%%%%%%%%%%%%%%%%%%%%
\begin{abstract}
This work presents a new enriched finite element method dedicated to the vibrations of axi- ally inhomogeneous Timoshenko beams. This method relies on the “half-hat” partition of unity and on an enrichment by solutions of the Timoshenko system corresponding to simple beams with a homogeneous or an exponentially-varying geometry. Moreover, the efficiency of the enrichment is considerably increased by introducing a new formulation based on a local rescaling of the Timoshenko problem, that accounts for the inhomogeneity of the beam. Val- idations using analytical solutions and comparisons with the classical high-order polynomial FEM, conduced for several inhomogeneous beams, show the efficiency of this approach in the time-harmonic domain. In particular low error levels are obtained over ranges of frequencies varying from a factor of one to thirty using fixed coarse meshes. Possible extensions to the research of natural frequencies of beams and to simulations of transient wave propagation are highlighted.
\end{abstract}

\begin{keyword}
  Wave-based finite element method \sep Timoshenko beams \sep vibrations \sep high frequency
\end{keyword}
%%%%%%%%%%%%%%%%%%%%%%%%%%%%%%%%%%%%%%%%%%%%%%%%%%%%%%%%%%%%%%%%%%%%%%%%

\maketitle


%
\section{Introduction}


Test de bibliographie :
 

%
\section{Notations and available exact solutions to the Timoshenko system}


\section{Enriched finite elements}
 


\section{A new formulation based on a local rescaling}

\vglue 2mm
...  /  ...
\vglue 2mm

%*********************************************
\section{Conclusions}



\section*{Acknowledgement}


%*********************************************
%*********************************************
\appendix

\section{Other enriched spaces}



%%%%%%%%%% BIBLIO %%%%%%%%%%%%%%%%%%%%%%%%%%%%%%%%%%%%%%%%%%%%%%%%%%%%%%%%
\bibliographystyle{abbrv}
\bibliography{Exercice_zotero2}


\end{document}